An $n$ is broken exactly when it has a prime factor $p\equiv1\pmod4$. If it has none, $n^2\mid x^2+y^2$ forces $n\mid x,y$: for odd primes $3\pmod4$, a square sum divisible by the prime forces both coordinates divisible, and repeat; for powers of two, repeated reduction modulo four does the same. Thus colours agree. Conversely for each relevant prime $p=5,13,17,29,37,41$, use a representation $p=u^2+v^2$ (respectively $(2,1),(3,2),(4,1),(5,2),(6,1),(5,4)$). The integer vector $(\dfrac{n}{p}) (u^2-v^2,2uv)$ has length $n$, but its coordinate sum is not divisible by $n$: if $p$ divided that sum, its square and the zero square-sum modulo $p$ would force one of the nonzero coordinates to vanish, impossible. Thus the broken values are $5,10,13,15,17,20,25,26,29,30,34,35,37,39,40,41,45,50$, summing to $\boxed{511}$.\begin{asy} import graph;import geometry;usepackage("amsmath");size(7cm);for(int x=0;x<=4;++x)for(int y=0;y<=5;++y)dot((x,y),gray);draw((0,0)--(3,0)--(3,4)--cycle);label("3",(1.5,0),(0,-1));label("4",(3,2),(1,0));label("5",(1.5,2),(-1,1));label("$(0,0)$",(0,0),(-1,-1));label("$(3,4)$",(3,4),(1,1));\end{asy}
